\scanpage{67}
\mspage{\handwritten{- 2 -}}

p.9, li. 2 du bas : qu'est-ce que désigne $e$ ?

p.10, li.7 : ``biéquidimensionnel (EGA $\unsure{0|O}_{\mathrm{IV}}$ 14.3.3).''

p.10, li. du bas : donner peut-être l'exemple de la fonction de dimension définie par un complexe dualisant (\handwritten{[}H\handwritten{]} V 7.1).

p.12, li. 7 du bas : ``telle que $G'$ soit ... soit ...''

p.13, li.11 du bas : ``continu pour la topologie ...''

p.14, li.12 : supprimer ``inversibles de $A$ de la forme''.

p.16; li. 3 : ``l'application composée ...''

p.16, li.9 : supprimer ``tel que $\dim(\operatorname{Supp}(\underline{F}))=j$''. Dire ce qu'on entend par``polynôme numérique''.

p.17 : pour la commutativit\unsure{\deleted{E}e|\deleted{e}E} du diagramme du bas, référer aux formules de changement de base (IV \hspace{3em}) ; la notation $\underline{s}_y$ devrait être remplacée par $s_{y_K}$ , où $y_K$ désigne l'image inverse de $y$ sur $X_K$ .

p.19, li.13 du bas : (IV 2.12.1).

p.20, li.6 : remplacer ``dans la catégorie dérivée'' par (IV 2.12.4).

p.22, li.8 : ``polynôme de Snapper (2.3.2) ...''

p.22, dans la formule (4.3.5), ne faut-il pas plutôt lire
\[
\handwritten{\alpha}_1+\ldots+\handwritten{\alpha}_d\handwritten{\leq}d\quad ?
\]

\deleted{p\unsure{.}22, li\unsure{.}6 du bas}

p.22, li.8 du bas : vider ``$=(\hspace{3em})$'' ?

p.24, li.8 du bas : ``$\dim(X)=\dim(Y)=d$''.

p.25, li.6 : vider la phrase ``On appelle équivalence numérique ...'', et rédiger ainsi (4.5.1) : ``... si les homomorphismes $\operatorname{Gr}_i(X)\handwritten{\longrightarrow}\underline{Z}$ (resp. $\operatorname{Gr}^i(X)\handwritten{\longrightarrow}\underline{Z}$) qu'ils définissent via le produit d'intersection (4.1.1) coïncident.''

p.25, li. 3 du bas : ``D\deleted{E}e 2.4.4 '' et non ``\unsure{\deleted{R}D|\deleted{D}R}e 2.4.3 ''.

Suggestion aux \S\S\ 2 et 4 : ajouter en remarque (ou exercice) que
% archive sheet number: 66
